
\def\ChapterCode{ALS}
\chapter
[\CO{[AED]} Defibrillation mit halbautomatischen Geräten]
{\CC{[AED]}\\Defibrillation mit halbautomatischen Geräten}
% \AasRsN{RS VIII}


\ChapterIntro
{%
Die Kardiopulmonale Reanimation -- bestehend aus
Herdruckmassage, Beatmung und Elektrotherapie mittels eines
Defibrillators -- ist eine der Kernaufgaben eines Rettungs-
oder notfallmedizinischen Dienstes. Die korrekte und
routinierte Durchführung der Reanimation hat großen Einfluss
auf das Überleben des Patienten bzw. seine Lebensqualität
nach der Entlassung aus der Spitalsbehandlung.

Dieses Kapitel behandelt die Basisreanimation mittels
eines halbautomatischen Defibrillators basierend auf den
Richtlinien des European Resuscitation Council. 
}{%
\begin{description}
  \item[Maintainer:] --
  \item[Autoren:] Diverse
  \item[Reviewer:] --
  \item[Version:] {\DocVersion} ({\DocVersionInfo})
  \item[SHA1:] {\DocGitCommit}
\end{description}

{\TxTeilkapitelAASS}
}

\SectionUnderConstruction{}



% \clearpage
\begin{figure}[H]
\begin{WidePar}
\Captionof{table}{Basic Life Support mit AED.\label{tab:algorithmus-bls-aed}}
\index{Basic Life Support}\index{BLS|see{Basic Life Support}}
\includegraphics[width=\linewidth]{./extern/BLS_mit_AED.pdf}
\end{WidePar}
\end{figure}

% \clearpage
\begin{figure}[H]
\begin{WidePar}
\Captionof{table}{Pediatric Life Support.\label{tab:algorithmus-pls}}
\includegraphics[width=\linewidth]{./extern/PLS.pdf}\index{Pediatric Life
Support}\index{BLS|see{Pediatric Life Support}}
\end{WidePar}
\end{figure}
\dictum[Dickie, in: Samuel Shem: \emph{House of God.}
\cite{HouseOfGodDe}]{Alles in Ordnung, auch wenn Sie in Panik geraten sind und
sich beschissen fühlen.
Ich kenne das. Es ist furchtbar, und es wird nicht das
letzte Mal sein.
Aber vergessen Sie nicht, was Sie gesehen haben. Regel Nr.
3: Bei Herzstillstand zuerst den \EE{eigenen} Puls fühlen.
}
\section{Anmerkungen}

\Text{Herzdruckmassage: sooft wie möglich!}{%
Die HDM muss so rasch beginnend und so unterbrechungsfrei wie
möglich sein.

Nach Abgabe des Schocks muss die Herzdruckmassage sofort
weitergeführt werden! Auch wenn die Defibrillation
erfolgreich war, d.\,h. ein ordnungsgemäßer Herzrhythmus
wiederhergestellt wurde, kann es einige Zeit dauern, bis
wieder ein ausreichender Kreislauf einsetzt
\cite{Sandroni2008,Skhirtladze2010}.

\Fazit{Die Unterbrechung der HDM durch den AED/Defibrillator
soll insgesamt nicht mehr als \E{5\Sek*} betragen!}
}{}{}{}{}

\Text{Qualität der Herzdruckmassage}{% 
Die richtige Durchführung der Herzdruckmassage hat
wesentlichen Einfluss auf das Behandlungsergebnis bzw. das
Überleben.

Die \EE{Kompressionsrate} muss mindestens \EE{100\PerMin*}
betragen. Die \EE{Drucktiefe} soll (beim Erwachsenen)
\EE{mindestens 5\CM* und maximal 6\CM*} betragen. Damit sich
das Herz mit Blut füllen kann ist eine komplette
\EE{Entlastung} wichtig: Während der Entspannungsphase darf
der Helfer nicht auf dem Brustkorb des Patienten lehnen.
\cite[p.1283]{Erc2010Sec2}
}{}{}{}{}

\Text{Infektionen}{% 
Die Ansteckungsgefahr mit Krankheiten im Rahmen einer
Reanimation sind sehr gering. Es wurden noch keine
Infektionen mit schwerwiegenden Erkrankungen wie HIV oder
Hepatitis C während einer Basisreanimation beschrieben.
\cite[p.1284]{Erc2010Sec2} }{}{}{}{}

\Text{Probleme bei der HDM während des Defi-Ladevorganges}{% 
\Ca{Der folgende Kommentar betrifft nur Helfer, welche nach
den ERC~2010-Leitlinien arbeiten!} 
In den ERC-Leitlinien von
2010 wird erstmals empfohlen, die Herzdruckmassage fortzusetzen,
\EE{während der Defibrillator sich auflädt}, d.\,h. nach der
Analyse wird weiter massiert, bis der Defi bereit zur
Schockabgabe ist. Dies soll die HDM-Pausen verringern. Bei
vielen AED-Geräten kann es jedoch zu
Problemen kommen, sie interpretieren die
reanimationsbedingten Störungen der EKG-Ableitung als
Kontaktverlust bzw. Bewegungsartefakt und brechen den
Ladevorgang ab.\footnote{Davon betroffen sind derzeit z.\,B. die
Geräte der Schiller Fred{\texttrademark}-Serie, der
Defigard{\texttrademark} 1002 und der
LifePak{\texttrademark} 500. Bezüglich anderer Geräte liegen
der Redaktion derzeit noch keine Informationen vor. Es muss
jedoch davon ausgegangen werden, dass auch andere Modelle
ein derartiges Verhalten zeigen können. \E{(Angaben ohne
Gewähr. Stand: 2011-01-19)}} Dadurch wird die Schockabgabe
verhindert. Daher kann diese Methode nicht generell
empfohlen werden, sondern darf nur durchgeführt werden, wenn
der Hersteller des Gerätes dies freigegeben hat.
\Commentary[Herzmassage während AED-Ladezyklus]{Lt.
vorläufiger Auskunft der Distributoren ist dies momentan
zumindest mit den AED-Geräten der Fa. Schiller
(Fred{\texttrademark}, Defigard{\texttrademark} 1002) nicht
möglich, da auch während der Aufladephase das EKG über die
Defi-Elektroden überwacht wird. Die elektrische Aktivität
durch Thoraxkompressionen wird dann als Bewegungsartefakt
interpretiert und es wird eine interne Sicherheitsentladung
eingeleitet. Für die Modelle der Fred-Reihe läuft momentan
eine Evaluierung bezüglich eines Softwareupdates, es wird
daher aller Voraussicht nach in Zukunft eine Lösung geben.
 
Für die Geräte der Corpuls{\texttrademark}-Reihe liegen
uns derzeit noch keine Informationen vor.  
Der LifePak{\texttrademark} 500 von Medtronic verhält sich
wie der Defigard{\texttrademark} 1002. Bezüglich der Typen
LifePak{\texttrademark} 1000, 12 und 15 liegen ebanfalls
noch keine Informationen vor. 
Es muss
 davon ausgegangen werden, dass auch andere Geräte
ein derartiges Verhalten zeigen. \E{(Angaben ohne
Gewähr. Stand: 2011-01-19)}}

\Cave{Manche Defibrillatoren brechen den Ladevorgang ab,
wenn währenddessen eine Herzdruckmassage durchgeführt wird.}

\Fazit{Während des Ladevorganges des Defibrillators darf nur
bei Verwendung \E{dafür} vorgesehener Geräte eine Herzdruckmassage
durchgeführt werden.} 
}{}{}{}{}

\ChapterEnd